\par\Needspace{6\baselineskip}\medskip
\hypertarget{native-087-article068}{}\pdfbookmark[2]{68.}{article-068}
\begin{center}68.\end{center}
Solutio problematis inversi, puta determinatio longitudinis et latitudinis ex ascensione recta et declinatione, eidem triangulo sphaerico superstruitur: formulae itaque supra traditae huic fini accomodabuntur per solam permutationem ipsius $b$ cum $\delta$, ipsiusque $l$ cum $-\alpha$. Etiam has formulas, propter usum frequentem, hic apposuisse haud pigebit:

Secundum methodum art. 66 habemus
\[
\begin{aligned}
\sin(45^0-\tfrac12 b)\sin\tfrac12(E-l)&=\cos(45^0+\tfrac12\alpha)\sin(45^0-\tfrac12(\epsilon+\delta))\\[.5em]
\sin(45^0-\tfrac12 b)\cos\tfrac12(E-l)&=\sin(45^0+\tfrac12\alpha)\cos(45^0-\tfrac12(\epsilon-\delta))\\[.5em]
\cos(45^0-\tfrac12 b)\sin\tfrac12(E+l)&=\sin(45^0+\tfrac12\alpha)\sin(45^0-\tfrac12(\epsilon-\delta))\\[.5em]
\cos(45^0-\tfrac12 b)\cos\tfrac12(E+l)&=\cos(45^0+\tfrac12\alpha)\cos(45^0-\tfrac12(\epsilon+\delta)).
\end{aligned}
\]

Contra ad instar methodi alterius art. 67 determinabimus angulum \SourceLeaf{088}auxiliarem $\zeta$ per aequationem
\[
\operatorname{tang}\zeta=\frac{\operatorname{tang}\delta}{\sin\alpha},
\]
\noindent eritque
\[
\begin{aligned}
\operatorname{tang}l&=\frac{\cos(\zeta-\epsilon)\operatorname{tang}\alpha}{\cos\zeta}\\[.5em]
\operatorname{tang}b&=\sin l\operatorname{tang}(\zeta-\epsilon).
\end{aligned}
\]
\noindent Ad calculi confirmationem adiungi poterit
\[
\cos b=\frac{\cos\delta\cos\alpha}{\cos l}=\frac{\cos(\zeta-\epsilon)\cos\delta\sin\alpha}{\cos\zeta\sin l}\cdot
\]
\noindent Pro determinatione ipsius $E$ inservient perinde ut in art. praec. aequationes
\[
\begin{aligned}
\cos E&=\frac{\sin\epsilon\cos\alpha}{\cos b}=\frac{\sin\epsilon\cos l}{\cos\delta}\\[.5em]
\cos b\cos\delta\sin E&=\cos\epsilon-\sin b\sin\delta.
\end{aligned}
\]

Variationes differentiales ipsarum $l$, $b$ hisce formulis exhibebuntur:
\[
\begin{aligned}
\mathrm{d}l&=\frac{\sin E\cos\delta}{\cos b}\,\mathrm{d}\alpha+\frac{\cos E}{\cos b}\,\mathrm{d}\delta\\[.5em]
\mathrm{d}b&=-\cos E\cos\delta\,\mathrm{d}\alpha+\sin E\,\mathrm{d}\delta.
\end{aligned}
\]
